\documentclass[11pt,a4paper]{article}

% ================================================================
% DADOS EDITAVEIS
% Substitua os campos abaixo pelos dados do aluno e pelas medicoes.
% Depois de preencher tambem os resultados calculados, troque
% \dadoscompletosfalse por \dadoscompletostrue.
% ================================================================
\newif\ifdadoscompletos
\dadoscompletostrue

\newcommand{\aluno}{Rodrigo Malato Navegantes}
\newcommand{\matricula}{2023038175}
\newcommand{\turma}{EA1}
\newcommand{\professor}{Luiz Machado}
\newcommand{\dataexperimento}{11 de agosto de 2026}

% Medicoes: sem vento / com vento
\newcommand{\Vsem}{32{,}00}
\newcommand{\Vcom}{32{,}00}
\newcommand{\Isem}{0{,}394}
\newcommand{\Icom}{0{,}394}
\newcommand{\Tssem}{86{,}7}
\newcommand{\Tscom}{34{,}5}
\newcommand{\Tinfsem}{24{,}0}
\newcommand{\Tinfcom}{24{,}0}
\newcommand{\velcom}{3{,}00}

% Resultados: sem vento / com vento
\newcommand{\qsem}{12{,}61}
\newcommand{\qcom}{12{,}61}
\newcommand{\DTsem}{62{,}7}
\newcommand{\DTcom}{10{,}5}
\newcommand{\hexpsem}{12{,}92}
\newcommand{\hexpcom}{77{,}17}
\newcommand{\uhexpsem}{0{,}62}
\newcommand{\uhexpcom}{10{,}90}
\newcommand{\hradsem}{4{,}87}
\newcommand{\hradcom}{3{,}76}
\newcommand{\hconvsem}{7{,}37}
\newcommand{\hconvcom}{39{,}73}
\newcommand{\hteosem}{12{,}23}
\newcommand{\hteocom}{43{,}49}
\newcommand{\desviosem}{5{,}6}
\newcommand{\desviocom}{77{,}4}

% ================================================================
% FORMATACAO
% ================================================================
\usepackage[T1]{fontenc}
\usepackage[utf8]{inputenc}
\usepackage[brazil]{babel}
\usepackage{lmodern}
\usepackage{microtype}
\usepackage[a4paper,margin=2.45cm]{geometry}
\usepackage{amsmath,amssymb}
\usepackage{graphicx}
\usepackage{booktabs,tabularx,array}
\usepackage{siunitx}
\usepackage{xcolor}
\usepackage{caption}
\usepackage{float}
\usepackage{enumitem}
\usepackage{titlesec}
\usepackage{fancyhdr}
\usepackage{listings}
\usepackage{hyperref}

\definecolor{azul}{HTML}{164A73}
\definecolor{cinza}{HTML}{555555}
\definecolor{fundonota}{HTML}{F3F6F8}

\sisetup{
  locale = DE,
  per-mode = symbol,
  exponent-product = \cdot,
  output-decimal-marker = {,},
  separate-uncertainty = true
}

\hypersetup{
  colorlinks=true,
  linkcolor=azul,
  urlcolor=azul,
  citecolor=azul,
  pdftitle={Transferência de calor por convecção e radiação combinadas em uma barra cilíndrica},
  pdfauthor={\aluno}
}

\titleformat{\section}
  {\large\bfseries\color{azul}}
  {\thesection}{0.65em}{}
\titleformat{\subsection}
  {\normalsize\bfseries\color{azul}}
  {\thesubsection}{0.65em}{}
\titlespacing*{\section}{0pt}{1.4em}{0.55em}
\titlespacing*{\subsection}{0pt}{1.0em}{0.35em}

\pagestyle{fancy}
\fancyhf{}
\lhead{EMA-103 --- Laboratório de Térmica}
\rhead{Prática 1}
\cfoot{\thepage}
\renewcommand{\headrulewidth}{0.4pt}
\setlength{\headheight}{14pt}

\setlength{\parindent}{1.15cm}
\setlength{\parskip}{0.2em}
\renewcommand{\arraystretch}{1.16}
\captionsetup{font=small,labelfont=bf}

\lstdefinestyle{pythonema}{
  language=Python,
  basicstyle=\ttfamily\scriptsize,
  keywordstyle=\color{azul}\bfseries,
  commentstyle=\color{cinza}\itshape,
  stringstyle=\color{teal!70!black},
  numbers=left,
  numberstyle=\tiny\color{cinza},
  numbersep=7pt,
  frame=single,
  rulecolor=\color{black!25},
  backgroundcolor=\color{fundonota},
  breaklines=true,
  breakatwhitespace=false,
  showstringspaces=false,
  keepspaces=true,
  columns=fullflexible,
  tabsize=4,
  captionpos=b,
  literate=
    {á}{{\'a}}1 {à}{{\`a}}1 {â}{{\^a}}1 {ã}{{\~a}}1
    {é}{{\'e}}1 {ê}{{\^e}}1 {í}{{\'i}}1
    {ó}{{\'o}}1 {ô}{{\^o}}1 {õ}{{\~o}}1
    {ú}{{\'u}}1 {ç}{{\c c}}1
    {Á}{{\'A}}1 {É}{{\'E}}1 {Í}{{\'I}}1
    {Ó}{{\'O}}1 {Ú}{{\'U}}1 {Ç}{{\c C}}1
    {²}{{$^2$}}1 {±}{{$\pm$}}1
}

\newcommand{\naoinformado}{\textemdash}
\newcommand{\nota}[1]{%
  \begin{center}
    \colorbox{fundonota}{\parbox{0.92\linewidth}{\small #1}}
  \end{center}
}

\begin{document}

\begin{center}
  {\small UNIVERSIDADE FEDERAL DE MINAS GERAIS\\
  Escola de Engenharia --- Engenharia Aeroespacial\\
  EMA-103 --- Laboratório de Térmica}

  \vspace{1.1cm}

  {\LARGE\bfseries\color{azul}
  Transferência de calor por convecção e radiação combinadas em uma barra cilíndrica}

  \vspace{0.9cm}

  \begin{tabular}{@{}rl@{}}
    \textbf{Aluno(a):} & \aluno \\
    \textbf{Matrícula:} & \matricula \\
    \textbf{Turma:} & \turma \\
    \textbf{Professor(a):} & \professor \\
    \textbf{Data da prática:} & \dataexperimento
  \end{tabular}
\end{center}

\vspace{0.65cm}
\hrule

\section{Objetivo}

Determinar experimentalmente o coeficiente combinado de transferência de calor por convecção e radiação de uma barra cilíndrica de aço, aquecida internamente por resistores elétricos e exposta ao ar ambiente. Foram consideradas duas condições: convecção natural, sem acionamento do soprador, e convecção forçada, com escoamento transversal de ar. Os valores experimentais devem ser comparados às estimativas obtidas por correlações da literatura, levando em conta as incertezas de medição.

\section{Fundamentação teórica}

Em regime permanente, admitindo desprezíveis as perdas elétricas e a condução de calor pelos apoios, a potência elétrica fornecida aos resistores é igual à taxa de calor dissipada pela barra:

\begin{equation}
  \dot Q = V_{AB} I.
  \label{eq:potencia}
\end{equation}

Desprezando as áreas das extremidades, a superfície lateral de troca térmica é

\begin{equation}
  S = \pi d L,
  \label{eq:area}
\end{equation}

e o coeficiente experimental combinado é calculado por

\begin{equation}
  h_{\mathrm{exp}} =
  \frac{\dot Q}{S\left(T_s-T_\infty\right)}.
  \label{eq:hexp}
\end{equation}

O coeficiente linearizado de radiação entre a barra e as superfícies vizinhas, adotadas à temperatura do ambiente, é

\begin{equation}
  h_{\mathrm{rad}} = \varepsilon\sigma
  \left(T_s^2+T_\infty^2\right)
  \left(T_s+T_\infty\right),
  \label{eq:hrad}
\end{equation}

em que as temperaturas são absolutas, $\varepsilon=0{,}60$ é a emissividade da barra e $\sigma=\SI{5.67e-8}{\watt\per\square\metre\per\kelvin\tothe{4}}$ é a constante de Stefan--Boltzmann.

Para a condição sem vento, o número de Nusselt médio para convecção natural sobre um cilindro horizontal pode ser estimado pela correlação de Churchill e Chu \cite{incropera}:

\begin{equation}
  \overline{Nu}_d =
  \left[
  0{,}60+
  \frac{0{,}387\,Ra_d^{1/6}}
  {\left[1+\left(0{,}559/Pr\right)^{9/16}\right]^{8/27}}
  \right]^2,
  \qquad
  h_{\mathrm{conv}}=\frac{\overline{Nu}_d k}{d}.
  \label{eq:natural}
\end{equation}

Nessa expressão, $Ra_d=Gr_dPr$, $Gr_d=g\beta(T_s-T_\infty)d^3/\nu^2$ e as propriedades do ar são avaliadas na temperatura de filme $T_f=(T_s+T_\infty)/2$.

Para a condição com vento, emprega-se a correlação de Churchill e Bernstein para escoamento transversal sobre um cilindro \cite{incropera}:

\begin{equation}
  \overline{Nu}_d = 0{,}3+
  \frac{0{,}62Re_d^{1/2}Pr^{1/3}}
  {\left[1+(0{,}4/Pr)^{2/3}\right]^{1/4}}
  \left[1+\left(\frac{Re_d}{282000}\right)^{5/8}\right]^{4/5},
  \label{eq:forcada}
\end{equation}

com $Re_d=\rho U d/\mu$ e $h_{\mathrm{conv}}=\overline{Nu}_d k/d$. Para o cálculo da densidade na saída do soprador, utilizou-se

\begin{equation}
  \rho=\frac{rP_{\mathrm{atm}}}{RT_\infty},
  \qquad
  r=1{,}2,\quad
  P_{\mathrm{atm}}=\SI{92.45}{\kilo\pascal},\quad
  R=\SI{0.287}{\kilo\joule\per\kilogram\per\kelvin}.
  \label{eq:densidade}
\end{equation}

Por fim, a estimativa teórica do coeficiente combinado é

\begin{equation}
  h_{\mathrm{teo}}=h_{\mathrm{conv}}+h_{\mathrm{rad}}.
  \label{eq:hteo}
\end{equation}

O tratamento numérico foi implementado em Python, com uso do CoolProp para a obtenção das propriedades termofísicas do ar. O código completo é apresentado no Apêndice~\ref{ap:codigo}.

\section{Materiais e método}

Foram utilizados uma barra cilíndrica de aço com resistores internos, uma fonte de tensão com indicação de tensão e corrente, um soprador de ar, um termopar tipo K conectado a um medidor de temperatura, um anemômetro, uma escala e um paquímetro. A montagem experimental é apresentada na Figura~\ref{fig:montagem}.

\begin{figure}[H]
  \centering
  \includegraphics[width=0.82\linewidth]{Imagens/montagem_labtermo_p1.png}
  \caption{Montagem experimental da Prática 1}
  \label{fig:montagem}
\end{figure}

Inicialmente foram medidos o comprimento e o diâmetro externo da barra. Em seguida, os resistores foram energizados e aguardou-se a aproximação do regime permanente. Registraram-se a tensão, a corrente elétrica, a temperatura superficial da barra e a temperatura ambiente. O procedimento foi realizado sem vento e repetido com o soprador acionado; nessa segunda condição, mediu-se também a velocidade do ar nas proximidades da barra.

A pressão atmosférica foi obtida no Mapa de Estações do Instituto Nacional de Meteorologia (INMET) \cite{inmet}, utilizando-se a estação meteorológica da Pampulha em 11 de agosto de 2026, às 17 h. Nesse horário, foram registradas pressão instantânea de \SI{924.5}{\hecto\pascal}, máxima de \SI{925.3}{\hecto\pascal} e mínima de \SI{924.5}{\hecto\pascal}. Nos cálculos foi adotada a pressão instantânea, equivalente a \SI{92.45}{\kilo\pascal}, por representar a condição no instante considerado.

\section{Resultados e discussão}

\subsection{Medições}

A Tabela~\ref{tab:medicoes} reúne as dimensões fornecidas no roteiro e os campos para as medições efetuadas na bancada.

\begin{table}[H]
  \centering
  \caption{Grandezas medidas e incertezas instrumentais.}
  \label{tab:medicoes}
  \small
  \begin{tabular}{lccc}
    \toprule
    \textbf{Grandeza} & \textbf{Sem vento} & \textbf{Com vento} & \textbf{Incerteza} \\
    \midrule
    $V_{AB}$ (\si{\volt}) & \Vsem & \Vcom & $\pm\,3\%$ \\
    $I$ (\si{\ampere}) & \Isem & \Icom & $\pm\,3\%$ \\
    $L$ (\si{\milli\metre}) & 195  & 195   & $\pm\,\SI{0.5}{\milli\metre}$ \\
    $d$ (\si{\milli\metre}) & 25,4 & 25,4 & $\pm\,\SI{0.1}{\milli\metre}$ \\
    $T_s$ (\si{\celsius}) & \Tssem & \Tscom & $\pm\,\SI{1}{\celsius}$ \\
    $T_\infty$ (\si{\celsius}) & \Tinfsem & \Tinfcom & $\pm\,\SI{1}{\celsius}$ \\
    $U$ (\si{\metre\per\second}) & \naoinformado & \velcom & $\pm\,5\%$ \\
    \bottomrule
  \end{tabular}
\end{table}

\subsection{Propagação das incertezas}

Admitindo medições independentes, a incerteza relativa da potência elétrica é

\begin{equation}
  \frac{u_{\dot Q}}{\dot Q}=
  \sqrt{\left(\frac{u_V}{V}\right)^2+
        \left(\frac{u_I}{I}\right)^2}
  =\sqrt{0{,}03^2+0{,}03^2}=4{,}24\%.
  \label{eq:uq}
\end{equation}

Para a área lateral, a combinação das incertezas de $d$ e $L$ resulta em

\begin{equation}
  \frac{u_S}{S}=
  \sqrt{\left(\frac{u_d}{d}\right)^2+
        \left(\frac{u_L}{L}\right)^2}
  =0{,}47\%.
  \label{eq:uarea}
\end{equation}

Como $\Delta T=T_s-T_\infty$, tem-se $u_{\Delta T}=\sqrt{u_{T_s}^2+u_{T_\infty}^2}=\SI{1.41}{\kelvin}$. Assim, a incerteza do coeficiente experimental é calculada por

\begin{equation}
  \frac{u_{h_{\mathrm{exp}}}}{h_{\mathrm{exp}}}=
  \sqrt{
    \left(\frac{u_V}{V}\right)^2+
    \left(\frac{u_I}{I}\right)^2+
    \left(\frac{u_d}{d}\right)^2+
    \left(\frac{u_L}{L}\right)^2+
    \left(\frac{u_{\Delta T}}{\Delta T}\right)^2}.
  \label{eq:uh}
\end{equation}

\subsection{Coeficientes de transferência de calor}

\begin{table}[H]
  \centering
  \caption{Resultados experimentais}
  \label{tab:resultados}
  \small
  \begin{tabular}{lcc}
    \toprule
    \textbf{Grandeza} & \textbf{Sem vento} & \textbf{Com vento} \\
    \midrule
    $\dot Q$ (\si{\watt}) & \qsem & \qcom \\
    $\Delta T$ (\si{\kelvin}) & \DTsem & \DTcom \\
    $h_{\mathrm{exp}}$ (\si{\watt\per\square\metre\per\kelvin}) & \hexpsem & \hexpcom \\
    $u_{h_{\mathrm{exp}}}$ (\si{\watt\per\square\metre\per\kelvin}) & \uhexpsem & \uhexpcom \\
    $h_{\mathrm{rad}}$ (\si{\watt\per\square\metre\per\kelvin}) & \hradsem & \hradcom \\
    $h_{\mathrm{conv}}$ (\si{\watt\per\square\metre\per\kelvin}) & \hconvsem & \hconvcom \\
    $h_{\mathrm{teo}}$ (\si{\watt\per\square\metre\per\kelvin}) & \hteosem & \hteocom \\
    Desvio relativo (\si{\percent}) & \desviosem & \desviocom \\
    \bottomrule
  \end{tabular}
\end{table}

O acionamento do soprador reduziu a diferença de temperatura de \SI{62.7}{\kelvin} para \SI{10.5}{\kelvin}, embora a potência elétrica tenha permanecido em aproximadamente \SI{12.61}{\watt}. Como consequência, o coeficiente experimental combinado aumentou de \SI{12.92}{\watt\per\square\metre\per\kelvin} para \SI{77.17}{\watt\per\square\metre\per\kelvin}. A contribuição radiativa correspondeu a cerca de 40\% do coeficiente teórico sem vento, mas a aproximadamente 9\% com vento, mostrando que a convecção se tornou dominante no escoamento forçado.

Sem vento, o desvio entre os coeficientes experimental e teórico foi de 5,6\%. Embora $h_{\mathrm{teo}}=\SI{12.23}{\watt\per\square\metre\per\kelvin}$ fique ligeiramente abaixo do intervalo experimental $h_{\mathrm{exp}}\pm u_{h_{\mathrm{exp}}}$, a diferença é pequena diante da dispersão típica das correlações de convecção natural. Com vento, entretanto, o desvio foi de 77,4\% e o valor teórico não pertence ao intervalo experimental. Essa diferença pode estar associada à baixa diferença de temperatura, que amplia a incerteza relativa, à medição local da velocidade do ar, à não uniformidade da temperatura superficial, ao resfriamento do próprio termopar pelo jato e às perdas de calor pelas extremidades e pelos apoios, contabilizadas como convecção na determinação de $h_{\mathrm{exp}}$.

\section{Conclusão}

Os coeficientes combinados obtidos foram
\begin{align*}
  h_{\mathrm{exp,sem}} &= \hexpsem\pm\uhexpsem,
  & h_{\mathrm{teo,sem}} &= \hteosem,\\
  h_{\mathrm{exp,com}} &= \hexpcom\pm\uhexpcom,
  & h_{\mathrm{teo,com}} &= \hteocom,
\end{align*}
em \si{\watt\per\square\metre\per\kelvin}. O soprador produziu aumento expressivo da transferência de calor e reduziu a temperatura superficial da barra. A condição sem vento apresentou boa concordância com a estimativa teórica. Na condição com vento, o coeficiente experimental foi substancialmente superior ao previsto pela correlação, indicando influência relevante das limitações experimentais e das hipóteses adotadas no modelo.

\begin{thebibliography}{9}

\bibitem{slides}
UFMG. \textit{EMA-103: Laboratório de Térmica --- Aulas 2 e 3: convecção e radiação em uma barra cilíndrica}. Material didático da disciplina, 2026.

\bibitem{incropera}
BERGMAN, T. L.; LAVINE, A. S.; INCROPERA, F. P.; DEWITT, D. P. \textit{Fundamentals of Heat and Mass Transfer}. 6. ed. Hoboken: Wiley, 2011.

\bibitem{inmet}
INSTITUTO NACIONAL DE METEOROLOGIA (INMET). \textit{Mapa de Estações}. Disponível em: \url{https://mapas.inmet.gov.br/}. Acesso em: 22 ago. 2026.

\end{thebibliography}

\appendix
\section{Código em Python}
\label{ap:codigo}

\lstinputlisting[
  style=pythonema,
  caption={Código utilizado no tratamento dos dados experimentais.},
  label={lst:codigo}
]{calculo_transferencia.py}

\end{document}
